\subsection{Nondiscrete valuations of residue characteristic zero}

The logarithmic argument can be extended beyond DVRs. The full
valuative log structure is essential here: replacing it by a fixed
finitely generated submonoid would, in general, change its Kummer
coverings and its arithmetic fundamental group.

\begin{claim}{B.LIT5}{Proved}{All-prime logarithmic comparison over a henselian rank-one valuation ring}
Let $R$ be a henselian rank-one valuation ring with fraction field $K$
and residue field $\kappa$ of characteristic zero. Let
$S\subset\mathbb P^1(K)$ be a finite set containing $0,1,\infty$,
and let $\mathscr X/R$ be its stable marked genus-zero model.
Endow $\operatorname{Spec}R$ with the canonical valuative log
structure $R\setminus\{0\}\to R$, and $\mathscr X$ with the
log structure of its vertical boundary and its marked sections.
Write $\mathscr X_s$ for its log closed fibre and
$U=\mathbb P^1_K\setminus S$. Then restriction gives equivalences
\[
 \operatorname{FEt}(U)
 \ \longleftarrow\
 \operatorname{KCov}(\mathscr X)
 \ \longrightarrow\
 \operatorname{KCov}(\mathscr X_s).                 \tag{B.LIT.1}
\]
Here $\operatorname{KCov}$ denotes all finite Kummer \'etale
coverings in the saturated, not necessarily fs, sense of
Achinger--H\"ubner--Lara--Stix. In that definition the underlying
scheme map is integral and need not be finite. The equivalences
are over the entire $BG_K$, and are compatible with finite valued
extensions and with morphisms between stable models obtained by
forgetting marks. Neither discreteness nor completeness is assumed.
\end{claim}
\begin{proof}
We state the exact source inputs. AHLS Lemma~5.3.2 says that a
microbial valuation ring of residue characteristic zero is a
filtered colimit of local excellent regular subalgebras
$\Lambda_i$ with local maps to $R$, containing a fixed
pseudouniformizer $\pi$, for which $V(\pi)_{\mathrm{red}}$ is
a strict normal crossings divisor. If $R$ is henselian the
$\Lambda_i$ may be chosen henselian. A nontrivial rank-one
valuation is microbial, because $R[1/\pi]=K$ for every
$0\ne\pi\in\mathfrak m_R$. Put
\[
 M=R\setminus\{0\},\qquad
 M_i=\Lambda_i\cap\Lambda_i[1/\pi]^\times.
\]
The paragraph following that lemma gives
$(M\to R)=\varinjlim_i(M_i\to\Lambda_i)$ as saturated prelog
rings; the associated log schemes $B_i$ are fs and log regular.
AHLS Proposition~4.2.6 identifies finite Kummer \'etale covering
categories with their filtered colimit for an affine charted
inverse system of qcqs saturated log schemes, in the sense of its
Definition~3.4.4. These are the precise continuity and approximation
theorems used below.\citep[Definition~3.4.4,
Proposition~4.2.6, Lemma~5.3.2]{LA_AHLS2024}

For completeness, the required stable family exists over $R$,
without a noetherian assumption. Hassett Theorem~2.1, with genus
zero and all weights one, gives a proper moduli stack and its
projective coarse scheme over $\mathbb Z$. Every stable labelled
genus-zero curve has trivial automorphism group, so in this case
the stack is the fine projective moduli scheme
$\overline M_{0,|S|}$. Its universal family and the valuative
criterion give $\mathscr X$. The finite-presentation moduli
problem over arbitrary bases is obtained from the initially
noetherian formulation by descent of its finite presentation.
The model and its generic identification are unique, and its base
change to an extension of valuation rings is the corresponding
stable model.\citep[Theorem~2.1]{LA_Hassett2003}

One can describe its finite chart data directly. For a finite set
of split marks, the finite tree of branching clusters is determined
by the valuations of pairwise differences. At a branching radius,
choose a difference $d\ne0$ of that valuation and blow up the ideal
$(x-a,d)$ in the appropriate smooth coordinate. Its charts are
$x-a=dz$ and $(x-a)w=d$. Repeat at the finitely many branching
radii, omitting or contracting components with fewer than three
special points. This is the stable model just described. Its node
charts have the form $xy=a$, for $0\ne a\in\mathfrak m_R$;
its marked charts have a smooth horizontal parameter. All nodes
and their branches are split, since the marks are split and their
dual graph is a labelled tree.

Over the base monoid $M$, the log charts are $M$, $M\oplus\mathbb N$
at a mark, and
\[
       M[e_1,e_2]/(e_1+e_2=a)
                                                        \tag{B.LIT.2}
\]
at a node. AHLS Example~2.4.5 proves that these node monoids are
saturated and that the maps from $M$ are smooth, vertical and
saturated for every saturated $M$ and every $a\in M$.
The standard vertical log structure agrees with these charts by
AHLS Proposition~5.3.1; the disjoint smooth horizontal sections
add the displayed $\mathbb N$ factors. The log-trivial open is $U$.

Descend the marked stable family, its finite chart cover, and the
chart gluings to a sufficiently large $\Lambda_i$. This is ordinary
finite-presentation descent; equivalently, descend the $R$-point
of the fine moduli scheme and pull back its universal family.
After increasing $i$, all the finitely many node parameters and
chart identities descend. For each such nonzero parameter $a$,
rank one gives $N$ with $\pi^N/a\in R$. Require $a$ and
$\pi^N/a$ to belong to $\Lambda_i$. Then $a\in M_i$; on the
regular snc base it is a unit times a monomial in the local
equations of $V(\pi)_{\mathrm{red}}$. Include also the inverses
of the finitely many unit coefficient differences which separate
sections and chart overlaps. All these requirements persist on a
cofinal subsystem.

Let $\mathscr X_i/B_i$ denote the resulting marked log curve.
The proper stable underlying scheme is an ordinary base change
of one fixed descended family. Its local log charts are the
versions of (B.LIT.2) and $M_i\oplus\mathbb N$ over $M_i$.
It is thus fs and log smooth over the log-regular $B_i$, hence
log regular by the log-regularity theorem recalled immediately
after AHLS Lemma~5.3.2. All integers are invertible: the
$\Lambda_i$ can be taken as $\mathbb Q$-algebras in its
characteristic-zero construction. Let $U_i$ be the log-trivial
open; its limit is $U$.

The node and marked monoid maps are saturated. Their saturated
log base changes consequently have the same underlying schemes
as ordinary base changes. The compatible finite strict \'etale
chart covers make $\{\mathscr X_i\}$ an affine charted system
in exactly AHLS Definition~3.4.4, with limit $\mathscr X$.
The same applies to the log closed fibres
$\mathscr X_{i,s_i}$. Indeed, all the maps $\Lambda_i\to R$
and the transition maps are local inclusions. Their induced
residue-field maps are injective: the inverse image of
$\mathfrak m_R$ in each ring is its maximal ideal. Those
residue fields therefore have colimit $\kappa$. Taking the same charts modulo their maximal
ideals therefore gives the log fibre $\mathscr X_s$ as their
saturated inverse limit.

Lepage Theorem~1.6 identifies finite Kummer \'etale covers of an
fs log-regular scheme with finite \'etale covers of its log-trivial
open, for primes invertible on that scheme. His Theorem~1.8
identifies finite Kummer \'etale covers of a connected fs log
scheme proper over a noetherian henselian base with those of its
log closed fibre. Apply both theorems to every $\mathscr X_i$,
with all primes. They give
\[
 \operatorname{FEt}(U_i)\ \longleftarrow\
 \operatorname{KCov}(\mathscr X_i)\ \longrightarrow\
 \operatorname{KCov}(\mathscr X_{i,s_i}).
\]
Ordinary finite \'etale continuity on the left and AHLS
Proposition~4.2.6 on the other two terms prove (B.LIT.1).
This is a limit argument using the full valuative monoid, not an
application of the fs theorems directly to $R$.
\citep[Theorems~1.6 and 1.8]{Lepage2013}

Apply the same argument to the log base, proper over itself.
It identifies the covering categories of that base and its log
closed point with $\operatorname{FEt}(K)$. Independently, AHLS
Corollary~4.4.9 says that the Kummer \'etale fundamental group of
the canonical log valuation ring is the maximal tame quotient of
$G_K$. Every finite extension is tame in residue characteristic
zero, so this is all of $G_K$. Since all the restriction functors
commute with pullback of base covers, (B.LIT.1) is over $BG_K$.

For a finite valued extension $L/K$, AHLS Proposition~4.4.3
identifies the canonical valuative log map as a finite Kummer
\'etale map in the stated integral sense. Stable models base
change as schemes, and their saturated node charts base change
as just checked. This yields compatibility with the actual field
extension. For a finite diagram of mark-forgetting morphisms,
descend the diagram and all its charts to a common stage; the
restriction squares commute at that stage and hence at the limit.
Equivalently, a transported cover with its specified restriction
has a contractible groupoid of extensions by full faithfulness.
The resulting comparison isomorphisms therefore satisfy their
cocycle identities. This proves the asserted functoriality.
\citep[Proposition~4.4.3, Corollary~4.4.9]{LA_AHLS2024}
\end{proof}

\begin{claim}{B.LIT6}{Refuted}{Reconstruction over all finite-transcendence-degree henselian fields of residue characteristic zero}
Every field $K$ admitting a nontrivial henselian rank-one valuation
of residue characteristic zero fails closed-point-preserving
reconstruction of finite-type $K$-schemes from their \'etale topoi
over $BG_K$. In particular this holds for every field of finite
transcendence degree over $\mathbb Q$ with a nontrivial henselian
valuation of residue characteristic zero, regardless of its rank,
discreteness, or transcendence defect.
\end{claim}
\begin{proof}
We first work in rank one. Choose $0\ne t\in\mathfrak m_K$ and
put $\lambda=1+t$. On $\mathbb P^1(\overline K)$ set
\[
 F(a)=\begin{cases}\lambda a,&v(a)>0,\\a,&v(a)\leq0,\end{cases}
 \qquad F(0)=0,\qquad F(\infty)=\infty.
                                                        \tag{B.LIT.3}
\]
Henselianity makes the extended valuation $G_K$-invariant.
Thus $F$ is an equivariant bijection, with inverse scaling the
same ball by $\lambda^{-1}$. As in \statementref{B.CF6},
\[
       \frac{F(a)-F(b)}{a-b}\in1+\mathfrak m_{\overline K}
       \quad(a\ne b).                                  \tag{B.LIT.4}
\]
When both points lie in the ball the ratio is $\lambda$; when
neither does it is one. In the remaining case
$v(a)>0\geq v(b)$, the change $ta$ has strictly greater
valuation than $a-b$. This proves (B.LIT.4) for an arbitrary
ordered rank-one value group. Hence every difference has the same
valuation and initial form before and after $F$.

Take a finite $G_K$-stable marked set containing $0,1,\infty$
and a finite Galois extension $L/K$ splitting it. The stable
models for $S$ and $F(S)$ have the same labelled cluster tree.
For a cluster with centre $a$ choose a scalar $d\in L^\times$
realizing its radius; use that \emph{same} scalar for the
transformed cluster with centre $F(a)$. Instead of the integral
power coordinates of the DVR proof, write $z=(x-a)/d$. If
$w=(x-b)/e$ is another component coordinate, the transition is
\[
                  z=\frac{b-a}{d}+\frac ed w.            \tag{B.LIT.5}
\]
Extract the coefficient of smaller valuation. Its first nonzero
normalized reduction is a nonzero constant, $w$, or $c+w$.
Equation (B.LIT.4) preserves this reduced factor and its value.
It therefore preserves both the ghost-monoid element and the
unit part of every logarithmic pullback. The smoothing parameter
$e/d$ can be chosen literally the same on the two sides.
Products, quotients and reciprocal coordinates give all node,
mark and infinity charts, including generators whose pullback
vanishes identically on a contracted tail. This proves
compatibility with every mark-forgetting contraction.

The labelled component maps lift uniquely once these ghost maps
and the log base are fixed. At a node a possible difference of
two lifts multiplies the branch generators by units $u_1,u_2$
with $u_1u_2=1$, $u_1x=x$, and $u_2y=y$. Thus
$u_1-1\in(x)\cap(y)=0$ in the reduced nodal ring. On the
smooth locus the vertical log structure comes from the fixed
base, while the horizontal generator is a non-zero-divisor.
Initial forms commute with finite extensions of valued fields,
so the comparisons commute with those extensions; their unique
labelled description and equivariance of $F$ give Galois descent.
At a mark $a$, the horizontal parameter changes by $\lambda$
or one. In particular $K(F(a))=K(a)$ as actual subfields,
the cusp inertia is preserved, and the quotient $G_{K(a)}$ is
identified by the identity. This is the same local Kummer-chart
calculation as in \statementref{B.CF6}; principal units have all
their roots by Hensel's lemma, since every integer is invertible.

Apply \statementref{B.LIT5} to these split models and descend.
It gives compatible equivalences between the finite \'etale
cover categories of $\mathbb P^1_K\setminus S$ and
$\mathbb P^1_K\setminus F(S)$ over the entire $BG_K$.
They respect all inclusions of marked sets and the actual
arithmetic cusp residue extensions. We recall why this transports
the whole small \'etale site, as required in
\statementref{B.CF7}. A quasi-compact separated \'etale object
over $\mathbb A^1_K$ is an open subscheme of the normalization
in its generic finite algebra. Its data consist of that finite
algebra, its branching over finitely many closed points, and the
retained unramified points with their residue extensions.
The cover equivalences and cusp comparisons carry all these
data. Maps extend uniquely from their generic algebra maps,
with the condition on retained points preserved. Geometric
point surjectivity, which detects covering families, is preserved
as well. Such objects form a basis of the small \'etale site.
The compatibility already proved therefore yields a pinned
topos self-equivalence over $BG_K$, whose function on
$\overline K$-valued points is (B.LIT.3).

Residue characteristic zero implies $v(q)=0$ for every
$q\in\mathbb Q^\times$. Thus $F$ fixes all rational numbers
and moves $t$. No $K$-polynomial has that point function: a
polynomial fixing infinitely many rational numbers is the
identity. The constructed topos morphism is consequently not
induced by a $K$-scheme morphism. This proves failure of
reconstruction in rank one.

Finally let $\operatorname{trdeg}_{\mathbb Q}K=d<\infty$
and let the original valuation have residue characteristic zero.
It is trivial on $\mathbb Q$. If
$v(a_1),\ldots,v(a_m)$ are rationally independent, distinct
monomials in the $a_i$ have distinct valuations, so no nonzero
$\mathbb Q$-polynomial in them can vanish. Hence $m\leq d$:
the value group has finite rational rank. Any strict chain of
convex subgroups gives rationally independent elements chosen
in successive differences; the highest nonzero coefficient in
a rational relation cannot vanish modulo the preceding convex
subgroup. The ordered group therefore has finite rank and a
maximal proper convex subgroup. Its quotient is a nonzero
rank-one ordered group. Coarsen by this subgroup.
ADH Lemma~3.4.2 states that a valued field is henselian if and
only if both a coarsening and its induced residue valuation are
henselian. The coarsening is therefore henselian; its residue
characteristic is still zero because $\mathbb Q^\times$
remains in its unit group. The rank-one result applies to the
same field $K$.\citep[Lemma~3.4.2]{ADH2017}
\end{proof}

This extension removes the nondiscrete residue-characteristic-zero
cases from the unresolved portion of B8. It does not replace the
all-prime comparison by a tame quotient in mixed characteristic.
Fields admitting no nontrivial henselian residue-characteristic-zero
coarsening require a separate argument.
